\labsection{5}{Robotics: Sensor-to-Action Control}{sec:lab05-robotics}

\begin{labproject}{Build and evaluate an interpretable wall-following controller}
\textbf{Scenario.} A SCITOS G5 robot uses ultrasound sensors to follow a wall. You will derive a compact reactive controller from real sensor readings and test its robustness.\\
\textbf{Software.} Python 3.11+, pandas, NumPy, scikit-learn, Matplotlib.\\
\textbf{Goal.} Convert four distance readings into one of four motion commands and evaluate the controller on unseen data.

\kagglesource{https://www.kaggle.com/datasets/uciml/wall-following-robot}{Sensor Readings from a Wall-Following Robot}
\kaggleexample{https://www.kaggle.com/code/thehapyone/multiple-classifier-robot-movement-classification}{Multiple Classifier Robot Movement Classification}
\repofolder{lab05\_robotics/}

\begin{labmode}
Download the dataset manually and place \texttt{sensor\_readings\_4.csv} under \texttt{data/wall\_following/}. This file contains \texttt{SD\_front}, \texttt{SD\_left}, \texttt{SD\_right}, \texttt{SD\_back}, and the robot action label.
\end{labmode}

\textbf{Required run.}
\begin{lstlisting}[style=mlpython]
python lab05_robotics/src/main.py \
  --data data/wall_following/sensor_readings_4.csv \
  --seed 42
\end{lstlisting}

\begin{mlsteps}
  \item \textbf{Inspect sensor behaviour.} Report class counts and summary statistics for the four distance sensors. Save two plots showing how front and side distances differ across actions.
  \item \textbf{Create a fixed split.} Use a stratified 70/30 train-test split with seed 42. Save the row indices under \texttt{config/split.csv}.
  \item \textbf{Implement a rule controller.} Use a small set of thresholds on front/left/right distance to return one of the four dataset actions: Move-Forward, Slight-Right-Turn, Sharp-Right-Turn, or Slight-Left-Turn. Tune thresholds on the training set only.
  \item \textbf{Create a reference model.} Fit a shallow decision tree with maximum depth 4 using the same four sensor inputs. The tree is a compact comparison model, not the final controller.
  \item \textbf{Evaluate on the test set.} Report accuracy, macro F1, per-class recall, and confusion matrices for both approaches. Identify which action is hardest to predict.
  \item \textbf{Test robustness.} Add zero-mean Gaussian noise at 2\%, 5\%, and 10\% of each sensor's training-set standard deviation. Re-evaluate the rule controller and plot macro F1 versus noise level.
\end{mlsteps}

\begin{pitfall}
Do not tune controller thresholds after viewing test labels. If the rule is changed using test errors, create a new holdout set before reporting the final metric.
\end{pitfall}

\requiredoutput{\texttt{metrics.csv}, two confusion matrices, sensor plots, noise-robustness plot, and the saved split.}
\end{labproject}

\begin{verification}
The controller must return exactly one valid action for every test row. Confirm that no test labels are used by the threshold-tuning function.
\end{verification}

\begin{labdeliverable}
Push the completed lab and create tag \texttt{lab05-final}. The README must state the final rules in plain language, compare them with the decision tree, and explain one reason real robot performance may differ from offline classification accuracy.
\end{labdeliverable}
